% エントリ 7(2026-11-03): gamma_BH の相対論的な補正の感度
\section{[E7, 2026-11-03] 自己相互作用の率の相対論的な補正: 感度}\label{sec:e07}

E4・E5 のあと残った最大の理論の不定性は、$211\times211\to322\times\mathrm{BH}$ の率 $\gamma_\mathrm{BH}$ の相対論的な補正である。準平衡で回転を失う速さは $1/\gamma_\mathrm{BH}$ に比例するので(式 \eqref{eq:e03_qe})、補正 $C_\mathrm{BH}(\alpha)\equiv\gamma_\mathrm{R}/\gamma_\mathrm{NR}$ は回転を失う時間に $C_\mathrm{BH}$ 倍で効く。年齢の比 $R$ には $C_\mathrm{BH}(\alpha_\mathrm{1G})/C_\mathrm{BH}(\alpha_\mathrm{2G})$ 倍で効く。計算は \texttt{calc/e07\_cbh.py}、出力は \texttt{calc/out/e07.json}。

\subsection{補正の値}

\arxiv{2412.03655} は Kerr 時空の Green 関数で率を計算し、非相対論の値との比を図で示している(同論文の Fig.~3、$\tilde a=0.95$)。数値の表はないので、図から目で読んだ。
\begin{itemize}
\item $\gamma^{322\times\mathrm{BH}}_{211\times211}$ の補正は次のとおり。読み取りの誤差は 20\% 程度とみる。
  \begin{center}\small
  \begin{tabular}{@{}l*{11}{c}@{}}
  \toprule
  $\alpha$ & 0.03 & 0.05 & 0.075 & 0.10 & 0.125 & 0.15 & 0.175 & 0.20 & 0.25 & 0.30 & 0.35 \\
  $C_\mathrm{BH}$ & 1.0 & 1.3 & 1.7 & 2.5 & 3.5 & 4.4 & 5.2 & 5.2 & 3.5 & 2.2 & 1.4 \\
  \bottomrule
  \end{tabular}
  \end{center}
\item 無限遠への放出 $\gamma^{211\times\infty}_{322\times322}$ の比は、$\alpha\le0.2$ で 0.9--1.2(同図の左)。$\dot\chi$ には平方根で入るので落とした。
\item BH のスピンへの依存は図にない。$\tilde a=0.95$ の値を $\chi=0.5$--$0.8$ にそのまま使った。
\end{itemize}
したがってこのエントリは感度の検定であり、本計算では補正を自分で計算するか、\arxiv{2412.03655} のコードの出力と照合する。

\subsection{結果}

\begin{figure}[H]
\centering
% E7: gamma_BH の相対論的な補正を入れたときの年齢の比(mu = 0.71e-12 eV)。データは calc/e07_cbh.py
\begin{tikzpicture}
\begin{axis}[note, width=0.62\linewidth, height=0.42\linewidth, ymode=log, xmin=6, xmax=28, ymin=1, ymax=2e3,
  xlabel={1 世代目の質量 $M_1\ [M_\odot]$}, ylabel={$R=t_\mathrm{1G}/t_\mathrm{2G}$},
  legend style={font=\scriptsize, at={(0.98,0.98)}, anchor=north east, fill=white, fill opacity=0.85, text opacity=1, draw=none}, legend cell align=left]
  \addplot[cA, line width=0.9pt] table[x=M1, y=R_d05] {figs/data/e07_R_mu0.71.dat}; \addlegendentry{主要項、$\delta=0.05$}
  \addplot[cA, line width=0.9pt, dashed] table[x=M1, y=R_d18] {figs/data/e07_R_mu0.71.dat}; \addlegendentry{主要項、$\delta=0.18$}
  \addplot[cB, line width=0.9pt] table[x=M1, y=Rc_d05] {figs/data/e07_R_mu0.71.dat}; \addlegendentry{相対論的な $\gamma_\mathrm{BH}$、$\delta=0.05$}
  \addplot[cB, line width=0.9pt, dashed] table[x=M1, y=Rc_d18] {figs/data/e07_R_mu0.71.dat}; \addlegendentry{相対論的な $\gamma_\mathrm{BH}$、$\delta=0.18$}
\end{axis}
\end{tikzpicture}

\caption{$\mu=0.71\times10^{-12}\eV$ での年齢の比 $R$。青は主要項の $\gamma_\mathrm{BH}$(E4 と同じ)、橙は相対論的な補正を入れた場合。超放射の率はどちらも E4 の相対論的な率。}
\label{fig:e07}
\end{figure}

\begin{result}
\begin{itemize}
\item \textbf{GW241011 が許す $f$ の上限}($T_2=10^7\yr$): $\mu=0.60$、$0.71$、$0.89\times10^{-12}\eV$ で、補正なし $6.1$、$4.1$、$2.5\times10^{12}\GeV$ から、補正あり $8.9$、$6.8$、$4.9\times10^{12}\GeV$ になる。GW241011 の $\alpha$($0.09$--$0.13$)で $C_\mathrm{BH}\approx2$--$4$ だからである。
\item \textbf{年齢の比}($\mu=0.71\times10^{-12}\eV$、$\chi_0=0.7$、$\delta=0.05/0.18$):
  \begin{center}
  \begin{tabular}{@{}lcc@{}}
  \toprule
  $M_1$ & 補正なし & 補正あり \\
  \midrule
  $8\Msun$ & 967/326 & 431/145 \\
  $10\Msun$ & 309/102 & 157/52 \\
  $15\Msun$ & 37/11 & 26/7.6 \\
  $20\Msun$ & 7.6/1.7 & 7.8/1.8 \\
  \bottomrule
  \end{tabular}
  \end{center}
  軽い 1 世代目は $\alpha$ が小さく $C_\mathrm{BH}$ も小さいので、比は小さくなる。窓は開く側に動く。
\item \textbf{結論への影響}: 窓が開くには、8--15 $\Msun$ の 1 世代目が 2 世代目より 8--430 倍年をとっている必要がある(補正あり)。補正は比を最大 2 倍動かすが、問いの形は変えない。本計算の補正の精度は 2 倍より十分よくなければならない。
\end{itemize}
\end{result}
